\clearpage\pagestyle{otMain}\setlength{\otBlocEspacement}{2.5mm}\fontsize{11}{15}\selectfont\corrigesChapitre{2}{Additionner et soustraire}
\begin{blocLivre}{cor-2-1}\begin{corrige}{1}{Deux sommes posées}\demarche{Aligner les unités sous les unités.}
En alignant les rangs : $6\,478+525=\res{7\,003}$ et $38\,467+4\,850=\res{43\,317}$. L’estimation $38\,000+5\,000=43\,000$ confirme l’ordre de grandeur.
\end{corrige}\end{blocLivre}
\begin{blocLivre}{cor-2-2}\begin{corrige}{2}{Deux différences posées}\demarche{Vérifier chaque différence par une addition.}
$5\,204-1\,368=\res{3\,836}$ ; contrôle : $3\,836+1\,368=5\,204$. Avec les zéros : $70\,000-2\,460=\res{67\,540}$.
\end{corrige}\end{blocLivre}
\begin{blocLivre}{cor-2-3}\begin{corrige}{3}{Compléments utiles}\demarche{Chercher ce qui manque pour atteindre le repère.}
$64+36=100$, donc le complément est $\res{36}$. $380+620=1\,000$, donc le complément est $\res{620}$.
\end{corrige}\end{blocLivre}
\begin{blocLivre}{cor-2-4}\begin{corrige}{4}{Ajouter 9, 19, 29 de tête}\demarche{Passer par la dizaine supérieure.}
On ajoute la dizaine entière puis on retire 1 : $457+10-1=\res{466}$ ; $457+20-1=\res{476}$ ; $262+30-1=\res{291}$.
\end{corrige}\end{blocLivre}
\begin{blocLivre}{cor-2-5}\begin{corrige}{5}{Estimer avant de calculer}\demarche{Arrondir chaque terme au millier.}
$4\,980$ est proche de $5\,000$ et $3\,120$ de $3\,000$ : la somme est proche de $8\,000$. Le seul résultat vraisemblable est $\res{8\,100}$.
\end{corrige}\end{blocLivre}
\begin{blocLivre}{cor-2-6}\begin{corrige}{6}{Regrouper pour aller vite}\demarche{Chercher deux termes dont la somme est ronde.}
On regroupe $254+46=300$, puis $300+378=\res{678}$. L’ordre des termes d’une somme peut être changé.
\end{corrige}\end{blocLivre}
\begin{blocLivre}{cor-2-7}\begin{corrige}{7}{Une soustraction à trous}\demarche{Utiliser l’opération qui défait l’autre.}
Le premier nombre vaut $385+247=\res{632}$. Pour le second : $902-345=\res{557}$. Contrôles : $632-247=385$ et $902-557=345$.
\end{corrige}\end{blocLivre}
\begin{blocLivre}{cor-2-8}\begin{corrige}{8}{Le bon panneau}\demarche{Comparer les deux distances restantes.}
La distance restante a diminué de $68-29=\res{39}$ km. Le mot « plus loin » n’indique pas une addition : c’est l’écart entre deux restes.
\end{corrige}\end{blocLivre}
\begin{blocLivre}{cor-2-9}\begin{corrige}{9}{Avec un schéma en barres}\demarche{Dessiner le tout et les parties connues.}
Les deux parties vendues totalisent $512+218=730$ billets. La partie manquante du schéma vaut $850-730=\res{120}$ places libres. Contrôle : $730+120=850$.
\label{fig-cor-2-9}\figureGeo[corrige]{barres-9}
\end{corrige}\end{blocLivre}
\begin{blocLivre}{cor-2-10}\begin{corrige}{10}{Le compte est-il bon ?}\demarche{Toujours contrôler une différence par la somme.}
$3\,584+1\,587=5\,171$, qui n’est pas $4\,003$ : le résultat est faux. Le bon calcul donne $4\,003-1\,587=\res{2\,416}$ ; contrôle : $2\,416+1\,587=4\,003$.
\end{corrige}\end{blocLivre}
\begin{blocLivre}{cor-2-11}\begin{corrige}{11}{La bibliothèque}\demarche{Suivre l’histoire étape par étape.}
Après les prêts, il reste $12\,540-3\,275=9\,265$ livres. Après le don, les étagères portent $9\,265+1\,860=\res{11\,125}$ livres. L’estimation $12\,500-3\,300+1\,900\approx11\,100$ confirme la réponse.
\end{corrige}\end{blocLivre}
\begin{blocLivre}{cor-2-12}\begin{corrige}{12}{Qui a le plus ?}\demarche{Repérer qui possède le plus avant de calculer.}
Nina a plus de cartes : Marco en a $325-68=\res{257}$. Ensemble, ils ont $325+257=\res{582}$ cartes. Le mot « plus » décrivait la comparaison, pas l’opération.
\label{fig-cor-2-12}\figureGeo[corrige]{barres-12}
\end{corrige}\end{blocLivre}
\begin{blocLivre}{cor-2-13}\begin{corrige}{13}{La cagnotte de la classe}\demarche{Calculer d’abord ce qui est déjà réuni.}
Les recettes totalisent $480+655=1\,135$ euros. Il manque $1\,500-1\,135=\res{365}$ euros. Contrôle : $1\,135+365=1\,500$.
\end{corrige}\end{blocLivre}
\begin{blocLivre}{cor-2-14}\begin{corrige}{14}{Trois étapes de course}\demarche{Déduire chaque donnée manquante dans l’ordre.}
La deuxième étape mesure $712-56=656$ km. Les deux premières totalisent $712+656=1\,368$ km. La troisième mesure $2\,000-1\,368=\res{632}$ km. Contrôle : $712+656+632=2\,000$.
\end{corrige}\end{blocLivre}
\begin{blocLivre}{cor-2-15}\begin{corrige}{15}{Le vrai du faux}\demarche{Conserver l’écart en déplaçant les deux nombres.}
Ajouter 2 aux deux nombres ne change pas leur écart : $\res{\text{elle a raison}}$. $759-400=\res{359}$, donc $757-398=359$. Sur une droite, les deux nombres ont glissé ensemble de 2.
\end{corrige}\end{blocLivre}
